% ============================================================
% 00-abstract.tex — V13 Rewritten Abstract
% ============================================================
\begin{abstract}
University counseling spans relational curricula, cohort-specific fees, exemption policies, and narrative regulations, each requiring distinct evidence representations and computational treatments. Controlled comparisons against functionally matched single-agent alternatives remain limited. We present CTU-Chat, a supervisor-routed multi-agent system that assigns queries to four bounded specialists with scoped tool privileges and representation-specific evidence paths: Neo4j graph traversal for curricula, structured lookup with deterministic calculation for financial data, and hybrid document retrieval with neural reranking for regulations. We evaluate CTU-Chat through three experimental scenarios on a 100-query held-out benchmark: architecture comparison, five-component ablation, and semantic-neighbor tool stress testing up to 51 tools. In architecture comparison, CTU-Chat obtains an 8.0\,pp higher End-to-End Success point estimate (81.0\% vs.\ 73.0\%; paired 95\% CI $[0.0, +16.0]$\,pp; McNemar $p{=}0.096$) and 41.4\% fewer input tokens at the cost of higher latency ($\sim 2\times$). Under registry expansion to 51 tools, scoped specialization preserves 88.0\% E2E execution (vs.\ 78.7\% for Top-10 selection), without displaying statistically divergent degradation slopes. These findings suggest that domain-scoped specialization can support institutional counseling reliability under the evaluated conditions, with measurable coordination trade-offs.

\keywords{Multi-Agent Systems \and Supervisor Routing \and Architectural Ablation \and Heterogeneous Knowledge Allocation \and Retrieval-Augmented Generation \and University Counseling.}
\end{abstract}
